\chapter{Conclusion and Future Work}

The findings from this study illustrate that there is a communication bottleneck within pipeline-parallel reinforcement learning. In general, there is considerable overhead associated with repeated gradient exchange between devices during the process of training. As shown in Chapter~3, the amount of data transferred during gradient communication can be nearly 2.4$\times$ more than the amount of data needed to send raw observations directly to a centralized server.

In order to address these issues of high communication costs, we introduced gradient accumulation along with percentile-based sparsification. Gradients are collected at the downstream device, but only the largest 10 percent of gradients by magnitude are sent to the upstream device each minibatch. A second important finding related to our initial approach is that simple accumulation without some form of decay will degrade the overall quality of the policy to approximately 53 percent of baseline performance. This occurs due to the fact that once gradients accumulate over several iterations without being sent, they become stale and no longer represent the current optimization landscape. By introducing a buffer decay factor which decays exponentially with $\lambda = 0.99$, we can solve this problem. The decay effectively acts as a relevance filter, where only recently significant gradients are allowed to reach the transmission threshold. With a retention factor of $0.99^{320} \approx 0.04$ per iteration, gradients that do not exceed the percentile threshold within a single iteration are effectively discarded.

We also found that the application of the 90th percentile along with buffer decay resulted in a reduction in total cumulative data transfer by approximately 35 percent, while also achieving a mean evaluation return of 572.52 versus 534.63 for the uncompressed baseline. We verified these results across multiple random seeds and observed both configurations converge to comparable performance.

There are several directions this work could be extended in the future. One possibility is to use adaptive percentile thresholds that adjust during training based on gradient statistics, which could achieve higher compression in later stages of training when gradients begin to stabilize. Another direction would be to combine gradient sparsification with quantization techniques to further reduce the amount of data being communicated. Finally, testing on more complex environments with larger networks would help determine how well the method scales beyond the CNN architecture used in this work.
